\section{Introduction and Business Motivation}

\subsection{Problem}
Scanning a room into a 3D model is already a consumer feature on iPhone~Pro and iPad~Pro through Apple RoomPlan and
apps such as Polycam, all of which rely on the \emph{LiDAR sensor} for directly metric depth. Devices with LiDAR are,
however, a minority of the market: non-Pro iPhones and almost every Android phone have none. The alternative for these
devices is \emph{monocular depth estimation}---networks that predict depth from a single RGB image---which advanced
rapidly between 2020 and 2024 (MiDaS, Depth Anything) to the point of producing visually convincing results on single images.

The question a product team must answer before investing is: \textbf{if LiDAR is replaced by a monocular model while the
rest of the system stays the same, by how many centimetres is the reconstructed room wrong, and which use cases does that
still serve} (floor-area estimation, furniture placement, renovation measurement)? Monocular-depth research reports
per-image metrics (AbsRel, $\delta_1$) on benchmarks, which do not translate into ``how crooked is the wall.''
End-to-end reconstruction systems (Atlas, NeuralRecon, SimpleRecon) design their whole pipeline around their own model,
so the \emph{depth source} cannot be compared in isolation.

Some recent models claim to output \emph{metres} directly, without calibration (DA-v2 Metric-Indoor, ZoeDepth, Depth
Pro)---if that held, the question above would reduce to choosing a model. These models, however, are trained on data
that are not a phone's camera (DA-v2 Metric-Indoor on rendered Hypersim images), and inside a real room-reconstruction
pipeline the room comes out \SI{\rsMETRICCM}{cm} wrong (\S\ref{sec:exp1}). The direct remedy is \emph{transfer
learning}: keep training the model on depth from the target domain. The question is where those labels come from---the
\$20k laser scanner used to build research datasets, or the \emph{LiDAR already inside Pro phones}, which an app maker
can collect at scale. If the latter works, the minority of devices with LiDAR can be the teacher of the majority
without it.

\subsection{Research question}
\textbf{Main question:} is DA-v2 Metric-Indoor fine-tuned with depth from the iPad's LiDAR (no laser scanner at
training time) \emph{better than the same model pretrained}, when measured against the Faro laser scanner---the raw
ground truth ARKitScenes ships---both as a room mesh (centimetres on the wall) and as per-frame depth, on rooms never
used for training?
\textbf{Supporting questions:} (a) how much better is a laser teacher than a phone-LiDAR teacher; (b) where does the
fine-tuned model stand relative to using a sensor at run time (LiDAR, sparse points per frame); (c) where does the
remaining error come from?

\subsection{Approach}
This is a \emph{system-development} study. We build a classical room-reconstruction pipeline (per-frame depth
$\rightarrow$ scale alignment $\rightarrow$ TSDF voxel fusion $\rightarrow$ mesh) in which \textbf{the depth source is
the single swappable variable}, so ``pretrained'' and ``fine-tuned'' are two rows that differ only in a checkpoint
name, and every other depth source can be priced through the same pipeline as context:

\begin{table*}[t]
\centering\small
\begin{tabularx}{\textwidth}{@{}lX@{}}
\toprule
Depth source & Stands for \\
\midrule
\textbf{DA-v2 Metric-Indoor pretrained / fine-tuned with a LiDAR teacher (R3)} & \textbf{main question: a device with no sensor and no calibration} \\
Faro laser (\code{highres\_depth}) & upper bound of the pipeline: its own loss with perfect depth \\
ARKit LiDAR (\code{lowres\_depth}, $256\times192$) & what iPhone~Pro does today (and R3's teacher) \\
Monocular (Depth Anything V2, MiDaS) $\times$ scale strategy & ceiling and cost of supplying scale at run time instead of training \\
\bottomrule
\end{tabularx}
\end{table*}

Fine-tuning uses ARKitScenes rooms disjoint from the test rooms at the \code{visit\_id} level and trains only the DPT
neck and metric head. Evaluation uses Faro twice: a reference mesh fused from Faro depth on six test rooms (3D), and Faro
depth per frame directly on those rooms and on \rsVN{} more from the Validation fold (2D). Isolating \emph{ScaleAligner}
as a separate stage further lets us separate ``wrong shape'' from ``wrong scale'': \code{oracle\_frame} (fit per frame
against GT) is a model's ceiling, \code{per\_scene} is a one-off calibration, and \code{sparse\_points} (fit per frame
on a few hundred metric points) is what a system with a VIO tracker can do.

\subsection{Main results}
\begin{itemize}
\item \textbf{The LiDAR-teacher fine-tune (R3) beats the pretrained model in every test room:} Chamfer
\SI{\rsMETRICCM}{cm}$\rightarrow$\SI{\rsFTLIDARALLCM}{cm} ($\rsMAINGAIN\times$, \rsMAINWINS/\rsMAINN{} rooms), \fscore{}
\rsMETRICF$\rightarrow$\rsFTLIDARALLF, per-frame AbsRel against Faro \rsMETRICSIXABSREL$\rightarrow$\rsFTLIDARALLABSREL{}
---with no calibration and no sensor at run time (\S\ref{sec:main}).
\item \textbf{The same holds on \rsVN{} rooms that took part in nothing:} AbsRel \rsVMETRICABSREL$\rightarrow$\rsVFTALLABSREL,
better in \rsVWINS/\rsVN{} rooms (exact Wilcoxon $p=\rsVP$) (\S\ref{sec:main}).
\item \textbf{A laser teacher is not measurably better than a phone-LiDAR teacher:} \SI{\rsFTFAROCM}{cm} (Faro) /
\SI{\rsFTLIDARCM}{cm} (LiDAR, 10 rooms) / \SI{\rsFTLIDARALLCM}{cm} (LiDAR, 24 rooms), within the across-room std
(\S\ref{sec:finetune}).
\item An intrinsics-conditioned metric model (Depth Pro), given the real focal length and no fine-tuning, reaches only
\SI{\rsDEPTHPROCM}{cm}---knowing the camera does not replace in-domain labels (\S\ref{sec:finetune}).
\item \textbf{R3 does not reach a sensor:} the pipeline itself loses \SI{\rsGTCM \pm \rsGTSD}{cm}, iPad LiDAR
\SI{\rsLIDARCM \pm \rsLIDARSD}{cm}, monocular depth with the correct scale on every frame
\SI{\rsORACLECM \pm \rsORACLESD}{cm}, and a per-frame fit on $\sim$200 sparse points (VIO proxy) \SI{\rsSPARSECM}{cm}
(\S\ref{sec:exp1}).
\item \textbf{The remaining error is scale that swings per view, not shape:} one calibration per room gives
\SI{\rsSCENECM \pm \rsSCENESD}{cm} uniformly; the model's scale varies several-fold within one scan with the content of
the frame, not with time (\S\ref{sec:discussion}). Fine-tuning removes the domain bias; the per-frame factor still has
to come from metric data at run time.
\item A $13\times$ smaller model (DA-v2 Small) costs \SI{1.0}{cm} for a $6\times$ speed-up, and frame count governs
coverage, not accuracy ($\geq 0.8$~fps is needed).
\end{itemize}

\subsection{Contributions}
\begin{enumerate}
\item \textbf{A measured answer to ``does fine-tuning on a phone's own LiDAR really help a metric monocular model?''},
against a laser scanner on held-out rooms in 3D (6 rooms) and 2D (6 $+$ \rsVN{} rooms), with the checkpoints, the room
split and the room-selection scripts released (\S\ref{sec:ftsetup}, \S\ref{sec:main}).
\item A teacher ablation (Faro vs.\ LiDAR, 10 vs.\ 24 rooms) and an intrinsics-conditioned baseline under one training
recipe (\S\ref{sec:finetune}).
\item An open-source pipeline in which \code{DepthSource} and \code{ScaleAligner} are interfaces switched from config
without touching the orchestrator, so every table is ``one loop, one variable'' and a newly trained checkpoint becomes
a new row immediately (\S\ref{sec:design}).
\item ``Centimetres on the wall'' for Faro / LiDAR / monocular $\times$ scale strategy through one pipeline, and an
analysis showing that monocular error comes from per-view scale, not temporal drift---which explains both why
fine-tuning helps and why it is not enough (\S\ref{sec:exp1}, \S\ref{sec:discussion}).
\item Per-scan cost (time, frames, model size, LiDAR confidence weighting) for business decisions (\S\ref{sec:business}).
\end{enumerate}

\paragraph{What this paper does \emph{not} claim}
It does not claim to beat Atlas / NeuralRecon / SimpleRecon---they are not compared, and this pipeline is not designed to
compete. It does not claim that fine-tuning makes LiDAR unnecessary: R3 (\SI{\rsFTLIDARALLCM}{cm}) remains far from
LiDAR (\SI{\rsLIDARCM}{cm}) and from having metric points at run time (\SI{\rsSPARSECM}{cm}). It does not claim that a
LiDAR teacher equals or beats a laser teacher---R1 and R2 differ in both label and number of training views
(\S\ref{sec:ftsetup}) and by less than the across-room std. Only the head of one model is fine-tuned (frozen encoder,
ARKitScenes, iPad camera); every other model is used pretrained, and other phone cameras are not measured.
